% critique.tex — f168r, vv. 466–508 : commentaire critique
% Articles textuels tirés du commentaire de l'auteur (latex/Aristophane-Thesmo Harmattan 2023 *.tex),
% texte et notes conservés ; répartition détaillée dans repartition.md.
\cc{477}{ξύνοιδ᾽ ἐμαυτῇ πολλ᾽ ἐκεῖνο δ᾽ οὖν.}{Il manque un pied au vers 477 ; Bentley y suppléa en ajoutant <δεινὰ> πολλ᾽ pour introduire le superlatif δεινότατον du vers suivant. C'est une hypothèse (Dawes préféra πολλὰ δείν᾽). J'en propose une autre qui a l'avantage d'expliquer paléographiquement la lacune par un \italique{saut du même au }quasi-\italique{même} (il n'y a qu'une patte de mouche de différence entre un π et un γ) : \textbf{ΠΟΛΛ'} à la place de ΠΟΛΛΑ\textbf{Γ'ΑΛΛ'}.}
\cc{493}{μάλισθ᾽ ὑπό του ληκώμεθα.}{Le scribe ici encore a hésité entre \grec{ληκώμεθα}, qu'il écrivit probablement en premier, avant de gratter les trois premières lettres et de réécrire par dessus \grec{κινώμεθα}. N'étant pas un aficionado de la \italique{lectio difficilior}, il aura probablement préféré le verbe \grec{κινεῖν}, bien connu de lui, au fort rare \italique{ληκᾶν} (ou \grec{ληκᾶσθαι} dans la bouche d'une femme). Suidas (λ 434) ne s'y est pas trompé qui glose ληκώμεθα par κινώμεθα en citant ce passage d'Aristophane.}
